\chapter{Entorno experimental}
\label{cap7:entorno}
\section{Descripción de los datos}
Los datos utilizados en este trabajo provienen del dataset Milán \cite{milandataset}. Este dataset fue creado con la intención original de descubrir la relación, si existiese, entre los eventos meteorológicos y el tráfico de red en la ciudad de Milán, Italia. El dataset contiene datos tanto climatológicos como de SMS, llamadas telefónicas y tráfico perteneciente a navegación web. Los datos están organizados en mediciones realizadas cada hora que corresponden a una celda cuadrada de $235 m$ de lado dentro de un área de $553.66km^2$ de la ciudad de Milán \cite{milanogrid}. El dataset contiene mediciones correspondientes a los meses de noviembre y diciembre de 2013. Para este trabajo se han aislado y agregado agregado los datos relativos a SMS, llamadas telefónicas y tráfico web en muestras equiespaciadas una hora en el tiempo, desechando la información relativa a los eventos meteorológicos.

Las celdas utilizadas en este trabajo corresponden a una cuadrícula de 7x7 celdas centrada en la celda 5060. Para averiguar el alcance de las redes neuronales convolucionales, también se han realizado pruebas con una cuadrícula de 15x15 celdas centradas en la celda 5464. Esta última cuadrícula engloba completamente a la cuadrícula 7x7.

Aproximadamente, todas las series temporales de las celdas son bastante similares entre sí. Todas presentan un patrón subyacente claro y periódico con un periodo de más o menos 24 horas. Además, todas ellas también presentan un patrón periódico con un periodo de una semana, en el que los primeros 5 ciclos de 24 horas tienen bastante más amplitud que los últimos 2. Esto coincide con los 5 días laborables de la semana y los 2 días del fin de semana.

Respecto a las similitudes entre celdas, es difícil afirmar que una celda tenga una serie similar a otra celda contigua a lo largo de toda la cuadrícula. Por ejemplo, la celda 5357 es contigua a las celdas 5257, 5258 y 5358 pero su serie temporal, aunque presenta las dos componentes periódicas comunes a toda la cuadrícula, no es similar como se aprecia en la Figura  \ref{5357_5358_raw}. Por otro lado, la celda 4962 es bastante similar a todas sus celdas contiguas (4861, 4862, 4863, 4961, 4963, 5061, 5062, 5063) presentando todas ellas las dos componentes cíclicas comunes de forma muy marcada.

\begin{figure}[h!]
	\centering
	\begin{subfigure}{\linewidth}
		\includesvg[inkscapelatex=false,scale=0.6]{Celdas/trafico_celda_5357.svg}
	\end{subfigure}
	\begin{subfigure}{\linewidth}
		\includesvg[inkscapelatex=false,scale=0.6]{Celdas/trafico_celda_5358.svg}
	\end{subfigure}
	\caption{Datos contenidos en el \textit{dataset} Milán para las celdas 5357 (arriba) y 5358 (abajo) \cite{milandataset}.}
	\label{5357_5358_raw}
\end{figure}
\FloatBarrier



La mayoría de las celdas extienden el patrón semanal hasta el día 22 de diciembre de 2013, a partir del cual el tráfico de red disminuye considerablemente. Esto coincide con las vacaciones de Navidad cuando los nodos que generan mayor tráfico como negocios, oficinas, universidades... están. Al gunas celdas, como las celdas 5357 y 5362 presentan dicha disminución del tráfico mucho antes, alrededor del día 1 de diciembre de 2013.
%%%%%%%%%%%BREVE REPRESENTACIÓN DE LOS DATOS

A la vista de estos cambios en los patrones del dataset Milán, se ha decidido entrenar los modelos de redes neuronales con las primeras 400 horas del dataset y utilizar el resto de horas, aproximadamente 1000 horas, como datos de validación. Este reparto asegura que los modelos son entrenados con un patrón invariable mientras el dataset de validación presenta cambios en el patrón subyacente. Así será posible observar si los modelos con \textit{online learning} son capaces de adaptarse a los cambios tal y como se ha planteado en los objetivos de este trabajo.




\section{Arquitecturas empleadas}
Definiendo el término arquitectura como la combinación de capas de neuronas que describe a una red neuronal y modelo como el conjunto de valores concretos de los hiperparámetros de una red neuronal, para este trabajo se han utilizado 4 arquitecturas de redes neuronales. En primer lugar se ha utilizado una arquitectura formada por una única capa de 24 neuronas LSTM en medio de una capa de entrada y una capa completamente conectada de 3 salidas. Con esta arquitectura se ha obtenido el rendimiento base propuesto en \cite{SKIMAcasaseca},a mejorar en este trabajo. Los datos procesados por esta arquitectura pertenecen a las celdas 4259, 4456 y 5060 del dataset Milán.

Se ha utilizado también una arquitectura formadas por una capa de entrada, una capa de neuronas LSTM y una capa completamente conectada. De esta forma se ha obtenido el rendimiento de la arquitectura exclusivamente LSTM\footnote{A lo largo de este trabajo se hará referencia a esta arquitectura como LSTM.} con las celdas contiguas de la ciudad de Milán.

Para intentar extraer la correlación espacial entre las celdas, se han propuesto otras 3 arquitecturas, resumidas en la Figura \ref{fig:arquitecturas}, compuestas por una capa de entrada, 1, 2 y 3 capas de neuronas CNN\footnote{A lo largo de este trabajo se hará referencia a estas arquitecturas como CNN+LSTM, 2CNN+LSTM y 3CNN+LSTM respectivamente.} respectivamente, una capa \texttt{Flatten}, una capa de neuronas LSTM y una capa completamente conectada para dar formato a las salidas \cite{timedistributed}.



\newpage
\begin{figure}[htbp]
	\centering
	\begin{tikzpicture}[
		>=Stealth,
		font=\sffamily,
		% Estilos de bloques reutilizables
		base block/.style={
			minimum width=2.4cm, 
			minimum height=0.72cm, 
			rounded corners=3pt, 
			align=center, 
			line width=0.7pt
		},
		title block/.style={
			base block,
			font=\small\bfseries,
			minimum height=0.62cm,
			draw=gray!50,
			fill=gray!10
		},
		input block/.style={
			base block,
			draw=teal!80!black, 
			fill=teal!20, 
			font=\scriptsize
		},
		conv block/.style={
			base block,
			draw=teal!70!black, 
			fill=teal!15, 
			font=\scriptsize
		},
		conv small/.style={
			base block,
			draw=teal!70!black, 
			fill=teal!15, 
			minimum height=0.52cm,
			font=\scriptsize
		},
		flatten block/.style={
			base block,
			draw=cyan!70!black, 
			fill=cyan!20, 
			font=\scriptsize
		},
		lstm block/.style={
			base block,
			draw=orange!85!black, 
			fill=orange!25, 
			minimum height=1.1cm, 
			font=\scriptsize\bfseries
		},
		dense block/.style={
			base block,
			draw=green!60!black, 
			fill=green!20, 
			minimum height=0.72cm, 
			font=\scriptsize
		},
		output block/.style={
			base block,
			draw=gray!60, 
			fill=white, 
			font=\scriptsize
		},
		arrow line/.style={->, thick, draw=black!75},
		column box/.style={
			draw=black!70, 
			rounded corners=6pt, 
			line width=0.8pt
		},
		td box/.style={
			draw=teal!70!black, 
			dashed, 
			rounded corners=4pt, 
			line width=0.7pt, 
			inner sep=3.5pt
		}
		]
		
		% =========================================================================
		% DEFINICIÓN DE ALTURAS HORIZONTALES FIJAS (Y) RECALIBRADAS Y ESPACIADAS
		% =========================================================================
		% Y_title   =  0.00
		% Y_input   = -1.15
		% Y_conv1   = -2.45
		% Y_flatten = -5.00  <-- Margen limpio garantizado bajo TimeDistributed
		% Y_lstm    = -6.40
		% Y_dense   = -7.75
		% Y_output  = -8.85
		% =========================================================================
		
		% -------------------------------------------------------------------------
		% COLUMNA 1: LSTM
		% -------------------------------------------------------------------------
		\node[title block]  (c1_title) at (0.0,  0.00)  {LSTM};
		\node[input block]  (c1_in)    at (0.0, -1.15)  {Input\\$(X_t, \text{LB} \times 49)$};
		\node[lstm block]   (c1_lstm)  at (0.0, -6.40)  {LSTM Layer};
		\node[dense block]  (c1_dense) at (0.0, -7.75)  {Dense Layer};
		\node[output block] (c1_out)   at (0.0, -8.85)  {Output\\$(\hat{Y}_t)$};
		
		\draw[arrow line] (c1_in) -- (c1_lstm);
		\draw[arrow line] (c1_lstm) -- (c1_dense);
		\draw[arrow line] (c1_dense) -- (c1_out);
		
		% -------------------------------------------------------------------------
		% COLUMNA 2: CNN + LSTM
		% -------------------------------------------------------------------------
		\node[title block]   (c2_title) at (3.5,  0.00)  {CNN+LSTM};
		\node[input block]   (c2_in)    at (3.5, -1.15)  {Input\\$(X_t, \text{LB} \times 7 \times 7)$};
		\node[conv block]    (c2_conv1) at (3.5, -2.45)  {Conv2D Layer 1};
		\node[flatten block] (c2_flat)  at (3.5, -5.00)  {Flatten};
		\node[lstm block]    (c2_lstm)  at (3.5, -6.40)  {LSTM Layer};
		\node[dense block]   (c2_dense) at (3.5, -7.75)  {Dense Layer};
		\node[output block]  (c2_out)   at (3.5, -8.85)  {Output\\$(\hat{Y}_t)$};
		
		\draw[arrow line] (c2_in) -- (c2_conv1);
		\draw[arrow line] (c2_conv1) -- (c2_flat);
		\draw[arrow line] (c2_flat) -- (c2_lstm);
		\draw[arrow line] (c2_lstm) -- (c2_dense);
		\draw[arrow line] (c2_dense) -- (c2_out);
		
		\begin{scope}[on background layer]
			\node[td box, fit=(c2_conv1)] (c2_td) {};
		\end{scope}
		
		% -------------------------------------------------------------------------
		% COLUMNA 3: 2CNN + LSTM
		% -------------------------------------------------------------------------
		\node[title block]   (c3_title) at (7.0,  0.00)  {2CNN+LSTM};
		\node[input block]   (c3_in)    at (7.0, -1.15)  {Input\\$(X_t, \text{LB} \times 7 \times 7)$};
		\node[conv block]    (c3_conv1) at (7.0, -2.45)  {Conv2D Layer 1};
		\node[conv block]    (c3_conv2) at (7.0, -3.55)  {Conv2D Layer 2};
		\node[flatten block] (c3_flat)  at (7.0, -5.00)  {Flatten};
		\node[lstm block]    (c3_lstm)  at (7.0, -6.40)  {LSTM Layer};
		\node[dense block]   (c3_dense) at (7.0, -7.75)  {Dense Layer};
		\node[output block]  (c3_out)   at (7.0, -8.85)  {Output\\$(\hat{Y}_t)$};
		
		\draw[arrow line] (c3_in) -- (c3_conv1);
		\draw[arrow line] (c3_conv1) -- (c3_conv2);
		\draw[arrow line] (c3_conv2) -- (c3_flat);
		\draw[arrow line] (c3_flat) -- (c3_lstm);
		\draw[arrow line] (c3_lstm) -- (c3_dense);
		\draw[arrow line] (c3_dense) -- (c3_out);
		
		\begin{scope}[on background layer]
			\node[td box, fit=(c3_conv1) (c3_conv2)] (c3_td) {};
		\end{scope}
		
		% -------------------------------------------------------------------------
		% COLUMNA 4: 3CNN + LSTM
		% -------------------------------------------------------------------------
		\node[title block]   (c4_title) at (10.5,  0.00) {3CNN+LSTM};
		\node[input block]   (c4_in)    at (10.5, -1.15) {Input\\$(X_t, \text{LB} \times 7 \times 7)$};
		\node[conv block]    (c4_conv1) at (10.5, -2.45) {Conv2D Layer 1};
		\node[conv small]    (c4_conv2) at (10.5, -3.30) {Conv2D Layer 2};
		\node[conv small]    (c4_conv3) at (10.5, -4.00) {Conv2D Layer 3};
		\node[flatten block] (c4_flat)  at (10.5, -5.00) {Flatten};
		\node[lstm block]    (c4_lstm)  at (10.5, -6.40) {LSTM Layer};
		\node[dense block]   (c4_dense) at (10.5, -7.75) {Dense Layer};
		\node[output block]  (c4_out)   at (10.5, -8.85) {Output\\$(\hat{Y}_t)$};
		
		\draw[arrow line] (c4_in) -- (c4_conv1);
		\draw[arrow line] (c4_conv1) -- (c4_conv2);
		\draw[arrow line] (c4_conv2) -- (c4_conv3);
		\draw[arrow line] (c4_conv3) -- (c4_flat);
		\draw[arrow line] (c4_flat) -- (c4_lstm);
		\draw[arrow line] (c4_lstm) -- (c4_dense);
		\draw[arrow line] (c4_dense) -- (c4_out);
		
		\begin{scope}[on background layer]
			\node[td box, fit=(c4_conv1) (c4_conv3)] (c4_td) {};
		\end{scope}
		
		% =========================================================================
		% CAJAS EXTERIORES AMPLIADAS HORIZONTAL Y VERTICALMENTE
		% Ancho aumentado a 3.10 cm (centrado ±1.55 cm)
		% Alto extendido de y = +0.65 cm (arriba) a y = -9.55 cm (abajo)
		% =========================================================================
		\draw[column box] (-1.55,  0.65) rectangle (1.55, -9.55);
		\draw[column box] ( 1.95,  0.65) rectangle (5.05, -9.55);
		\draw[column box] ( 5.45,  0.65) rectangle (8.55, -9.55);
		\draw[column box] ( 8.95,  0.65) rectangle (12.05, -9.55);
		
	\end{tikzpicture}
	\caption{Comparación de las arquitecturas evaluadas: LSTM base, CNN+LSTM, 2CNN+LSTM y 3CNN+LSTM para una cuadrícula de 49 celdas. Elaboración propia.}
	\label{fig:arquitecturas}
\end{figure}
\FloatBarrier

Para poder evaluar el rendimiento de las arquitecturas mencionadas en las dos cuadrículas se modifica simplemente el tamaño de las capas convolucionales para que se ajusten a las dimensiones de los conjuntos de celdas. El hecho de utilizar dos conjuntos, de 49 y 225 celdas, no implica utilizar una arquitectura distinta para cada conjunto. La arquitectura refiere a la disposición de las diferentes capas de neuronas en el modelo, no del tamaño ni la cantidad de neuronas presentes en ellas.

Para aprovechar al máximo las librerías cuDNN se ha utilizado la función de activación \texttt{sigmoide} para las puertas de entrada, salida y olvido y la función \texttt{tanh} para la combinación de entradas y estado de las neuronas LSTM \cite{lstmactivation,nvidiacudnn}. Otras funciones de activación no son compatibles con los \textit{kernels} optimizados cuDNN por lo que el tiempo de ejecución aumenta exponencialmente. Aunque cabe la posibilidad de que otras funciones de activación consigan mejores resultados, un tiempo de ejecución excesivamente alto inutiliza la red neuronal.

\section{Procedimiento de entrenamiento}
Para la evaluación y comparación de los modelos, se ha ideado un procedimiento similar a la validación de modelos de aprendizaje automático descrita en el Capítulo \ref{cap4:redes}. Se ha entrenado los modelos con las 400 primeras muestras del dataset y se han realizado dos validaciones distintas en cada modelo. Una de las validaciones se ha realizado con las otras 1000 muestras del dataset sin utilizar entrenamiento incremental. La otra validación se ha realizado con las mismas 1000 muestras pero implementando entrenamiento incremental con ajuste de gradientes por norma euclídea del error normalizado \cite{SKIMAcasaseca} definida por la ecuación:

\begin{equation}
	\norm{\frac{\textbf{Y}_{t-1}-\hat{\textbf{Y}}_{t-1}}{\textbf{Y}_{t-1}}}^2
	\label{eq:euclideannorm}
\end{equation}

Y el error cuadrático como función de coste:

\begin{equation}
	E_{t-1}(\textbf{W}_{t-1})=(\textbf{Y}_{t-1}-\hat{\textbf{Y}}_{t-1})^2
\end{equation}

De esta forma, las ecuaciones \ref{eq7} y \ref{eq:euclideannorm} se combinan en:

\begin{equation}
	\textbf{W}_t=\textbf{W}_{t-1}-\eta\cdot\norm{\frac{\textbf{Y}_{t-1}-\hat{\textbf{Y}}_{t-1}}{\textbf{Y}_{t-1}}}^2 \cdot \nabla_\textbf{W}(\textbf{Y}_{t-1}-\hat{\textbf{Y}}_{t-1})^2
\end{equation}

De esta forma se da una mayor importancia a las predicciones que han generado errores mayores. Estos errores dan más información y permiten realizar una actualización de pesos más agresiva para aumentar la velocidad de convergencia hacia el mínimo de la función de coste.

\section{Hiperparámetros considerados}
Hay una selección variada de parámetros configurables a la hora de entrenar las redes neuronales propuestas. No todos los parámetros han sido considerados como configurables como, por ejemplo, el tamaño de los sets de entrenamiento, test estático y test incremental. El criterio de selección se basa en el impacto de los hiperparámetros en el rendimiento de las redes neuronales en el apartado de \textit{online learning}. Por esta razón se ha seleccionado los siguientes parámetros:

\begin{itemize}
	\item \textbf{Número de celdas:} indica si se está procesando una cuadrícula de 7x7 celdas o de 15x15 celdas.

	\item \textbf{Tamaño del \textit{kernel}:} indica el tamaño del campo de visión de las neuronas convolucionales. Está estrechamente relacionado con la cantidad de celdas utilizadas en el dataset. Tiene un gran impacto en todas las fases de entrenamiento y test.

	\item \textbf{Tamaño de ventana deslizante:} también denominado \texttt{look\_back}, además de ser el número de neuronas presentes en la capa LSTM, establece cuántas muestras de tráfico de red de cada celda son procesadas al mismo tiempo. Dichas muestras son tomadas como una misma secuencia de longitud \texttt{look\_back}. Tiene un gran impacto en todas las fases de entrenamiento y test.

	\item \textbf{Tamaño de \textit{batch}:} también denominado \texttt{batch\_size}, indica cuántas secuencias de longitud \texttt{look\_back} son procesadas por la red neuronal en un mismo momento. Tiene un gran impacto en todas las fases de entrenamiento y test.

	\item \textbf{\texttt{n\_epoch\_updates}:} a partir de aquí \texttt{updates} para abreviar, indica la cantidad de veces que la red neuronal procesa un mismo \textit{batch} en \textit{online learning}. Cada vez que lo procesa, ajusta sus pesos si el error supera el umbral de entrenamiento incremental.

	\item \textbf{Umbral para el entrenamiento incremental:} también llamado \texttt{threshold}, indica el error mínimo que la red neuronal ha de cometer durante el entrenamiento incremental para ajustar sus pesos. Un mayor \texttt{threshold} elimina el ruido en los datos para evitar ajustes de pesos innecesarios aunque puede provocar que la red neuronal ignore cambios pequeños en el patrón subyacente de los datos.

	\item \textbf{Número de filtros de las capas CNN:} también llamado \texttt{num\_filter1}, \texttt{num\_filter2} o \texttt{num\_filter3}, indica cuántas formas de patrones son utilizadas por las capas convolucionales para intentar componer el patrón subyacente de los datos. Tiene un gran impacto en todas las fases de entrenamiento y test.
\end{itemize}

Para la búsqueda de los parámetros óptimos se ha preferido la estrategia \textit{grid search} frente a \textit{random search}. Los hiperparámetros considerados tienen un gran impacto en el tiempo de procesado de las redes neuronales, que se intenta reducir en la medida de lo posible. Esto establece un límite en el número de filtros, \texttt{updates} y el tamaño de la ventana deslizante y los \textit{batchs}. Además, un umbral \texttt{threshold} demasiado alto inutiliza el entrenamiento incremental. De esta forma, aunque \textit{random search} tenga posibilidades de probar combinaciones con valores fuera de lo normal, es casi seguro que éstos desemboquen en combinaciones que impongan una carga computacional excesiva.


\section{Métricas de rendimiento}
Como métrica de rendimiento se ha utilizado la raíz normalizada del error cuadrático medio o NRMSE (\textit{Root Normalized Mean Squared Error}). Sin embargo, como cada celda tiene su propio patrón temporal subyacente, cada celda tiene un NRMSE distinto. Teniendo 49 o 225 celdas es imposible comparar todos sus NRMSE con los NRMSE de otras celdas con otra combinación de hiperparámetros. Para dicha comparación, como todas las celdas son consideradas igual de importantes, se ha usado la media aritmética del NRMSE de todas las celdas.

El NRMSE eleva al cuadrado las diferencias entre el valor real $\textbf{Y}_i$ y el valor predicho por la red neuronal $\hat{\textbf{y}}_i$ y las promedia. Finalmente calcula la raíz cuadrada del promedio y normaliza, en este trabajo, por la desviación típica de los valores reales. El NRMSE queda por tanto definido como:

\begin{equation*}
	NRMSE=\frac{\sqrt{\frac{1}{n}*\sum_{i=1}^{n}(\textbf{Y}_i-\hat{\textbf{Y}}_i)^2}}{\sigma(\textbf{Y})}
\end{equation*}

donde $n$ es el número de muestras del conjunto medido.

El NRMSE es una métrica adimensional independiente de la magnitud de los datos que representa la distancia entre dos puntos. Es ideal para este trabajo puesto que los datos utilizados presentan grandes fluctuaciones que distorsionarían el error cuadrático medio y su raíz. Un valor de 0 de NRMSE sería perfecto, sin errores cometidos en todo el set. Sin embargo, es irracional buscar un NRMSE=0. Teniendo en cuenta los resultados obtenidos en \cite{SKIMAcasaseca}, el NRMSE buscado es aquel NRMSE<0,1778 en el set de entrenamiento incremental.


\section{Análisis de Spearman}
Para estimar el impacto de cada hiperparámetro en el desempeño de cada arquitectura, se ha decidido optar por el análisis de Spearman. Es un método no paramétrico para la estimación de la dependencia estadística entre dos variables. A diferencia del análisis de Pearson, que busca relaciones estrictamente lineales, Spearman busca cualquier tipo de relación.

Como todo algoritmo estadístico, el análisis de Spearman comienza por la obtención de una muestra de una población de dos variables. En el caso de este TFG una de las variables siempre es la métrica de rendimiento y la otra es cada uno de los hiperparámetros considerados.

La muestra obtenida ha de ser representativa de la población. En este TFG se ha realizado un \textit{grid search} con más de 2500 combinaciones, lo que arroja el mismo número de valores de la métrica de rendimiento NRMSE. Es un volumen suficientemente grande como para considerarla altamente representativa de la población.

Ciertamente las combinaciones de hiperparámetros incluidas en el \textit{grid search} de este TFG no son aleatorias, lo que reduciría la representatividad de la muestra. Sin embargo, se toma como población todas aquellas combinaciones de hiperparámetros que puedan, bajo un punto de vista lógico, arrojar un rendimiento aceptable en cuanto a recursos de computación y precisión de las predicciones. El resto de combinaciones son absolutamente inservibles por lo que no tenerlas en cuenta no produce un desequilibrio en la representatividad. Por ejemplo, cualquier combinación de hiperparámetros basada en \texttt{num\_filter1}=1 nunca conseguirá una predicción cercana al valor real.

Una vez obtenidas las muestras, el análisis de Spearman busca patrones comunes en ambas, es decir, estima la correlación existente entre ellas. La correlación implica un movimiento en la misma dirección en los valores de ambas variables a lo largo de toda la muestra. Dicho movimiento puede realizarse en el mismo sentido o en sentidos opuestos. Existe una correlación si al modificar el valor de un hiperparámetro se produce una modificación equivalente en el valor del NRMSE.

Al ser un algoritmo no paramétrico, no es necesario conocer la distribución de las poblaciones analizadas y es posible comparar variables con rangos de valores muy dispares y con muchas anomalías sin distorsionar el resultado. Esto es posible debido al uso de niveles ordenados que reemplazan los posibles valores de las variables. Por ejemplo, si la métrica de rendimiento NRMSE toma los valores 0.3, 0.25 y 0.6742 son sustituidos por los niveles $NRMSE_1$, $NRMSE_2$ y $NRMSE_3$ respectivamente y los valores de un hiperparámetro: 1580, 40721 y 31.67 son sustituidos por los niveles $H_2$, $H_3$ e $H_1$.

El análisis de Spearman calcula la diferencia al cuadrado entre los niveles de la métrica de rendimiento y el hiperparámetro contra el que se compara en cada uno de los puntos conocidos. Finalmente realiza una sencilla operación matemática para obtener el valor final de la correlación entre ambas variables:

\begin{equation*}
	r_s=1-\frac{6*\sum d^2}{n(n^2-1)}
\end{equation*}


siendo $d$ las diferencias entre niveles y $n$ el número de puntos existentes. El valor del coeficiente de Spearman $r_s$ varía entre -1 para variables completamente correlacionadas de forma inversa y 1 para una correlación total directa, siendo 0 el valor que indica la ausencia de correlación \cite{spearmansource}. Si un hiperparámetro está correlacionado con la métrica de rendimiento significa que un cambio en el primero empuja la métrica hacia un cambio en el mismo sentido. 

Al realizar múltiples veces el análisis de Spearman cambiando el hperparámetro considerado es posible conocer qué hiperparámetros y en qué medida influyen en la precisión de las predicciones.

Debido a la disparidad en la distribución de las variables que se analizan en este trabajo, por ejemplo: el hiperparámetro \texttt{look\_back} toma valores entre 12 y 32 mientras la métrica NRMSE del modelo con \textit{online learning} toma valores entre 0 y 1; el análisis de Spearman es ideal al ser capaz de evitar que la diferencia entre los rangos de valores de las diferentes variables afecte al cálculo de la correlación entre ellas.

Finalmente, aunque la correlación no necesariamente implica causalidad, es un buen método para estimar el impacto de los hiperparámetros en la precisión de las arquitecturas estudiadas.




